\section{Kebenaran Parsial dan Kebenaran Total}

Terdapat perbedaan penting antara dua jenis kebenaran program yang harus dipahami. \logic{Kebenaran parsial} (\textit{partial correctness}) dan \logic{kebenaran total} (\textit{total correctness}) berbeda pada satu aspek krusial: apakah program dijamin berterminasi (berhenti) \cite{rosen2019}.

\subsection{Kebenaran Parsial}

\begin{definisi}[Kebenaran Parsial]
Program $S$ dikatakan \textbf{benar secara parsial} terhadap prekondisi $p$ dan postkondisi $q$, ditulis $\{p\}\ S\ \{q\}$, jika: setiap kali eksekusi $S$ dimulai dalam keadaan yang memenuhi $p$ \textbf{dan $S$ berterminasi}, maka keadaan akhir memenuhi $q$.
\end{definisi}

Kata kunci dari kebenaran parsial adalah ``\textit{dan $S$ berterminasi}'' --- kebenaran parsial tidak memberikan jaminan bahwa program akan berhenti. Sebuah program yang mengandung infinite loop secara vakum memenuhi kebenaran parsial: karena program tidak pernah berterminasi, postkondisi trivially terpenuhi (tidak ada keadaan akhir yang perlu diperiksa).

\subsection{Kebenaran Total}

\begin{definisi}[Kebenaran Total]
Program $S$ dikatakan \textbf{benar secara total} terhadap prekondisi $p$ dan postkondisi $q$, ditulis $[p]\ S\ [q]$, jika: setiap kali eksekusi $S$ dimulai dalam keadaan yang memenuhi $p$, maka $S$ \textbf{pasti berterminasi} dan keadaan akhir memenuhi $q$.
\end{definisi}

Kebenaran total = kebenaran parsial + terminasi. Untuk membuktikan kebenaran total, kita harus membuktikan \emph{dua} hal secara terpisah:
\begin{enumerate}
    \item Kebenaran parsial: $\{p\}\ S\ \{q\}$ (menggunakan Hoare Triple dan loop invariant).
    \item Terminasi: program $S$ pasti berhenti ketika dimulai dari keadaan yang memenuhi $p$.
\end{enumerate}

\subsection{Bukti Terminasi}

Untuk membuktikan bahwa perulangan \texttt{while $B$ do $S$} pasti berterminasi, digunakan \logic{fungsi pengurangan} (\textit{variant function} atau \textit{loop variant}):

\begin{definisi}[Fungsi Pengurangan]
\logic{Fungsi pengurangan} (loop variant) adalah fungsi $V$ yang memetakan keadaan program ke bilangan bulat non-negatif, dan memenuhi:
\begin{enumerate}
    \item $V \geq 0$ selama kondisi perulangan $B$ bernilai benar.
    \item $V$ \textbf{menurun secara ketat} pada setiap iterasi perulangan.
\end{enumerate}
Karena bilangan bulat non-negatif tidak dapat menurun tanpa batas, perulangan harus berhenti setelah jumlah iterasi yang berhingga.
\end{definisi}

\begin{contoh}
\textbf{Terminasi Penjumlahan Array}

Pada program penjumlahan array (section sebelumnya), fungsi pengurangan adalah:
\[ V = n - i \]

\textbf{Verifikasi}:
\begin{itemize}
    \item Selama $i < n$ (kondisi loop): $V = n - i > 0$. $\checkmark$
    \item Pada setiap iterasi: $i$ bertambah 1, sehingga $V$ berkurang 1. $\checkmark$
\end{itemize}

Karena $V$ selalu menurun dan bernilai non-negatif, loop pasti berterminasi setelah paling banyak $n$ iterasi. $\blacksquare$
\end{contoh}

\subsection{Ringkasan Hubungan}

\begin{center}
\renewcommand{\arraystretch}{1.3}
\begin{tabular}{|p{4cm}|p{4cm}|p{4.5cm}|}
\hline
\textbf{Aspek} & \textbf{Kebenaran Parsial} & \textbf{Kebenaran Total} \\ \hline
Notasi & $\{p\}\ S\ \{q\}$ & $[p]\ S\ [q]$ \\ \hline
Jaminan terminasi & Tidak & Ya \\ \hline
Jaminan postkondisi & Hanya jika berterminasi & Selalu (pasti berterminasi) \\ \hline
Yang dibuktikan & Loop invariant & Loop invariant + fungsi pengurangan \\ \hline
\end{tabular}
\end{center}
